\clearpage\corrigesChapitre{13}{Inertie et chute libre}
\corrigesSouples
\begin{blocLivre}{cor-13-1}\begin{corrige}{1}{Un livre immobile}
\schema[corrige]{inertie}
\begin{questions}
\item Oui, dans le référentiel terrestre assimilé galiléen, son vecteur vitesse reste nul.
\item Non : poids et action de la table sont présents et se compensent.
\end{questions}
\end{corrige}\end{blocLivre}
\begin{blocLivre}{cor-13-2}\begin{corrige}{2}{Un palet idéal}
\begin{questions}
\item Il suit un mouvement rectiligne uniforme.
\item Non. Une vitesse constante ne nécessite pas une force dans le sens du mouvement.
\end{questions}
\end{corrige}\end{blocLivre}
\begin{blocLivre}{cor-13-3}\begin{corrige}{3}{Une voiture qui freine}
\begin{questions}
\item Non : sa valeur change même si sa direction reste la même.
\item Non, d’après la contraposée du principe d’inertie dans le référentiel adopté.
\end{questions}
\end{corrige}\end{blocLivre}
\begin{blocLivre}{cor-13-4}\begin{corrige}{4}{Définir la chute libre}
\begin{questions}
\item Le poids de la balle.
\item Non, la chute libre peut commencer avec une vitesse initiale non nulle.
\end{questions}
\end{corrige}\end{blocLivre}
\begin{blocLivre}{cor-13-5}\begin{corrige}{5}{Forces compensées}
\begin{questions}
\item La somme est nulle.
\item Non : il peut être immobile ou en mouvement rectiligne uniforme, selon sa vitesse initiale.
\end{questions}
\end{corrige}\end{blocLivre}
\begin{blocLivre}{cor-13-6}\begin{corrige}{6}{Au sommet du lancer}
\schema[corrige]{chute}
\begin{questions}
\item La vitesse y est nulle à l’instant du retournement ; le poids demeure vertical vers le bas et non nul.
\item Le vecteur vitesse évolue avant et après cet instant, et le poids n’est compensé par aucune autre force.
\end{questions}
\end{corrige}\end{blocLivre}
\begin{blocLivre}{cor-13-7}\begin{corrige}{7}{Un cercle uniforme}
\schema[corrige]{mouvement-circulaire}
\begin{questions}
\item La direction change constamment, tandis que la valeur reste constante.
\item Les forces ne se compensent pas dans le référentiel galiléen considéré.
\end{questions}
\end{corrige}\end{blocLivre}
\begin{blocLivre}{cor-13-8}\begin{corrige}{8}{Une vitesse limite}
\begin{questions}
\item Oui : le vecteur vitesse est constant.
\item Non : le poids non nul doit être compensé par d’autres actions du liquide. Ce n’est pas une chute libre.
\end{questions}
\end{corrige}\end{blocLivre}
\begin{blocLivre}{cor-13-9}\begin{corrige}{9}{Interpréter des vecteurs}
\begin{questions}
\item La direction et le sens restent les mêmes ; la valeur augmente.
\item Non, le vecteur varie : les forces ne se compensent pas dans les conditions du principe.
\end{questions}
\end{corrige}\end{blocLivre}
\begin{blocLivre}{cor-13-10}\begin{corrige}{10}{Une balle montante}
\begin{questions}
\item Non : le contact a cessé, donc l’action de contact de la main a disparu.
\item Le poids reste vers le bas tandis que la vitesse est vers le haut ; la vitesse varie et la montée ralentit.
\end{questions}
\end{corrige}\end{blocLivre}
\begin{blocLivre}{cor-13-11}\begin{corrige}{11}{Choisir le système}
\begin{questions}
\item Oui, dans le modèle simple où ces deux actions sont les seules pertinentes.
\item Non : ces deux forces agissent sur le même système. La réciproque du poids est l’attraction de la personne sur la Terre.
\end{questions}
\end{corrige}\end{blocLivre}
\begin{blocLivre}{cor-13-12}\begin{corrige}{12}{Un diagnostic trop rapide}
\demarche{Séparer la valeur de la vitesse et son orientation.}
\begin{questions}
\item Non : connaître le sens de la vitesse et le caractère non nul de la somme ne suffit pas à choisir une augmentation de valeur.
\item La valeur peut augmenter ou diminuer, la direction peut changer, ou plusieurs de ces effets peuvent se combiner.
\end{questions}
\retiens{Une somme des forces non nulle ne signifie pas automatiquement une augmentation de la rapidité.}
\end{corrige}\end{blocLivre}

